\documentclass{cls/magnonics}
\usepackage[utf8]{inputenc}
\usepackage[T1]{fontenc}
\usepackage{geometry}
\usepackage{graphicx}
\usepackage{float}
\usepackage{hyperref}
\usepackage[english,indonesian]{babel}

\title{Optimasi Pemrosesan Template LaTeX Menggunakan Pendekatan Berbasis AI}
\author{Academic Support Agent}
\date{\today}

\begin{document}
\maketitle

\begin{abstract}

\end{abstract}

\chapter{PENDAHULUAN}


\chapter{TINJAUAN PUSTAKA}


% Original content (commented out):
%  \documentclass{cls/magnonics}
\usepackage[english]{babel}

%%%%%%%%%%%%% INTRODUCTION %%%%%%%%%%%%%%%%%%%%%%%%%%%%%%%%%%%%%%%%
% Make sure that you read first the following comments and that you abide by the rules indicated in the different sections. This is because the content of this file will be parsed by an algorithm to form the booklet. 
% We recommend that you perform the edition inside overleaf (https://www.overleaf.com/project): a collaborative cloud-based LaTeX editor used for writing scientific documents.
% References shall be inserted in the file mylib.bib using the bibtex format. One can generate automatically bibtex entries from the DOI-number via https://www.doi2bib.org/. 

%%%%%%%%%%%%% PREAMBLES %%%%%%%%%%%%%%%%%%%%%%%%%%%%%%%%%%%%%%%%%%
% PLEASE DO NOT MODIFY THE FOLLOWING LINES
\usepackage{array}
\usepackage{graphicx}
\usepackage[utf8]{inputenc}
\usepackage{csquotes}
\usepackage{silence}
\WarningsOff[everypage]
\usepackage{background}
\usepackage{bm,amsmath,color,xcolor}
\usepackage{multirow,tabularx}
\usepackage{longtable}
\usepackage[backend=biber, maxbibnames=5, minbibnames=3, style=phys]{biblatex}

%%%%%%%%%%%%% TITLE AND AUTHORS %%%%%%%%%%%%%%%%%%%%%%%%%%%%%%%%
%% Please fill-in below between the {} the main information about the presenting author and her/his abstract. Use small letters (capitalization is managed automatically). It requires to protect between {} subscript and superscript text larger than 1 character: see example below. 

\def \mytitle  {title of the abstract for magnonics 2023 \textit{e.g.} spin-waves in Y\textsubscript{3}Fe\textsubscript{5}O{\textsubscript{12}}}
\def \absttype {invited / oral / poster}
\def \speaker  {presenting author}
\def \myemail  {presenting.author@email.edu}
\def \myaffil  {department name, institution 1, country}

%% PLEASE DO NOT MODIFY THE FOLLOWING LINES
\title{\titlecap{\mytitle}}
\author[1,*]{\titlecap{\speaker}}
\affil[1]{\titlecap{\myaffil}}
\corrauthor{\myemail}

%%%%%%%%%%%%% TITLE AND AUTHORS %%%%%%%%%%%%%%%%%%%%%%%%%%%%%%%%
% List here co-authors or alternate affiliations
\author[1,2]{\titlecap{second author}}
\author[2]{\titlecap{third author}}
\affil[2]{\titlecap{department name, institution 2, country}}

%%%%%%%%%%%%% BIBLIOGRAPHY %%%%%%%%%%%%%%%%%%%%%%%%%%%%%%%%%%%%%
% Bibliography is located in the mylib.bib file
\addbibresource{mylib.bib}

%%%%%%%%%%%%% DOCUMENT %%%%%%%%%%%%%%%%%%%%%%%%%%%%%%%%%%%%%%%%%
\graphicspath{{./figs/}}
\begin{document}
\maketitle

Enter your abstract here. Make sure that you read first carefully the comments inside this document and that you abide by the rules indicated. This is because the content of this file will be parsed by an algorithm to form the booklet. Note that the abstract should not be longer than 1 page in length. We recommend that you perform the edition inside the overleaf environment (\texttt{https://www.overleaf.com/project}): a collaborative cloud-based editor allowing free writing of \LaTeX\, documents.

It is a good idea to start your abstract with some background and motivation, then describe your approach and summarize the key results. We recommend adding an eye-catching figure that can be referred inside the text through its associated label \textit{e.g.} Figure.~\ref{myfig1} or  Eq.(\ref{eq:1}).

%%%%%%%%%%%%% EQUATIONS %%%%%%%%%%%%%%%%%%%%%%%%%%%%%%%%%%%%%%%%

\begin{equation}
    \frac {\partial \mathbf{M} }{\partial t} = - \gamma \mathbf{M} \times \mathbf{H}_\text{eff} + \alpha  \left (  \mathbf{M} \times \frac {\partial \mathbf{M} }{\partial t} \right ) \label{eq:1}
\end{equation}

%%%%%%%%%%%%% FIGURES %%%%%%%%%%%%%%%%%%%%%%%%%%%%%%%%%%%%%%%%%%

\begin{figure}[htpb!]
\begin{center}
% Figures shall be uploaded in the >figs folder
\includegraphics[scale=0.15]{vortex_mx.png}
\caption{Magnetic vortex state.}
\label{myfig1}
\end{center}
\end{figure}

References  shall be inserted in the \texttt{mylib.bib}-file using the bibtex format \footnote{Bibtex entries can be generated automatically from the DOI via web-tools such as \texttt{https://www.doi2bib.org/}}.  Cite them in the text using commands \cite{Li2018,Brataas2020} or \cite{Klein1993}.  They will appear automatically together at the end of the abstract sorted by quotation order. Do not try to modify the quotation style.

Final check: make sure that the total length of the abstract is no more than 1 page and that the file compiles without error before uploading it as a zip project inside the conference website. To save the project, click on the home symbol in the upper left corner of the overleaf editor, and download the project as a zip file. The project should be renamed using your correspondingauthor entry.
%%%%%%%%%%%%% REFERENCES %%%%%%%%%%%%%%%%%%%%%%%%%%%%%%%%%%%%%%%
\printbibliography[heading=bibliography]

%%%%%%%%%%%%% ADD WATERMARK %%%%%%%%%%%%%%%%%%%%%%%%%%%%%%%%%%%%
\backgroundsetup{
  position=current page.east,
  angle=-0,
  nodeanchor=east,
  vshift=+70mm,
  hshift=-17mm,
  opacity=1,
  color=black,
  scale=2,
  contents=\absttype
}
\end{document}

\end{document}